\documentclass[10pt,a4paper]{amsart}
\usepackage[margin=19mm]{geometry}
\usepackage{amsmath,amssymb,mathtools}
\usepackage[scr=rsfs,cal=euler]{mathalfa}
\usepackage{xfrac}
\usepackage[unicode,hidelinks,bookmarks=false]{hyperref}
\newcommand{\ie}{i.e.\ }
\newcommand{\cf}{cf.\ }
\newcommand{\resp}{respectively\ }
\newcommand{\Subsection}{\subsection}
\providecommand{\nobreakdash}{\nobreak\hbox{-}}
\setlength{\emergencystretch}{2em}
\multlinegap0pt
\usepackage{bm}
\usepackage{bbm}
\usepackage{dsfont}
\usepackage{xspace}
\usepackage{cancel}
\usepackage{mathtools}
\usepackage{amsthm}
\makeatletter
\@ifundefined{teo}{\newtheorem{teo}{Teo}}{}
\@ifundefined{lema}{\newtheorem{lema}[teo]{Lemma}}{}
\makeatother
\def\Dom{\text{\normalfont Dom}}
\newcommand{\A}{\mathcal{L}_a^s}
\newcommand{\N}{{\mathbb N}}
\newcommand{\R}{{\mathbb R}}
\title[An eigenvalue problem for a nonlocal quasilinear anisotropic]{An eigenvalue problem for a nonlocal quasilinear anisotropic equation in fractional Orlicz Sobolev spaces without the $\Delta\_2$--condition}
\author{Julian Fernandez Bonder, Martin Guzman, Juan F. Spedaletti}
\date{Source: arXiv:2603.20095v2 (CC BY 4.0)}
\begin{document}
\maketitle

\noindent\emph{Source and licence.} An eigenvalue problem for a nonlocal quasilinear anisotropic equation in fractional Orlicz Sobolev spaces without the $\Delta_2$--condition. Version arXiv:2603.20095v2 (CC BY 4.0), distributed under the Creative Commons Attribution 4.0 International licence, as recorded in the arXiv licence metadata for this exact version and shown on the abstract page https://arxiv.org/abs/2603.20095v2. Full rights notice: licenses/arxiv-2603.20095-CC-BY-4.0.txt.\par\smallskip

\noindent\textit{Editorial scope.} Counted unit: the Teo whose source label is \texttt{pseudo} in the article (source0.tex lines 485--504, source label \texttt{pseudo}), delivered together with the article's own complete original proof of it (source0.tex lines 593--676, one \texttt{proof} environment of 84 lines). No mathematics is written, repaired or re-derived: statement and proof are the article's own text line for line. The proof is detached from the statement in the source: the article states the result at lines 485--504 and prints its proof at lines 593--676, and the proof environment's own header names the statement, so the association is the article's own. Required context retained uncounted: MC (lines 413--416); paraEulerLagrange (lines 548--557). Every definition, display and label of those lines is retained unchanged. Omitted from this package and not redistributed: 1--412, 417--484, 505--547, 558--592, 677--1050. Nothing inside a retained excerpt is omitted or altered: every retained body is delivered exactly as the article's lines are written, so no omission receipt of any kind accompanies this package. Citations: the retained excerpts carry no citation command at all, so no bibliography entry of the article is transcribed and no reference is redistributed. Reference adaptation: none was needed and none was made; every reference inside the delivered unit resolves inside it, so no unresolved reference or citation is produced. Printed slips kept literal: no printed word, symbol, hypothesis or number of the counted statement or of its proof was altered. Layout and class: the article's own class is \texttt{amsart} with packages including geometry, graphicx, amsmath, hyperref; the wrapper is the lane's pinned \texttt{amsart} editorial wrapper at 10pt on A4, which restates the theorem-like environments the retained text needs and adds the article's own macro definitions that the retained text needs (\texttt{\textbackslash A}, \texttt{\textbackslash Dom}, \texttt{\textbackslash N}, \texttt{\textbackslash R}), with extra wrapper packages (bm, bbm, dsfont, xspace, cancel, mathtools). The delivered numbering is the unit's own. Assets: no retained span contains a figure environment, a graphics-inclusion command, a PDF-page inclusion, a table or a TikZ picture; no image file is copied into this package, the package declares no asset dependency and no image file is redistributed with it. Licence: version arXiv:2603.20095v2 of this article is distributed under the Creative Commons Attribution 4.0 International licence, as recorded in the arXiv licence metadata for this exact version and shown on the abstract page https://arxiv.org/abs/2603.20095v2. A full rights notice is delivered at licenses/arxiv-2603.20095-CC-BY-4.0.txt. Characters: every retained excerpt is pure ASCII, so the delivered engine, XeLaTeX with its default Latin Modern text font, reports no missing character.
\begin{equation}\label{MC}
(a(x,y,\xi)-a(x,y,\xi'))(\xi-\xi')\geq 0.
%\tag{MC}
\end{equation}

\begin{lema}\label{paraEulerLagrange}
If there exists $u \in W_0^sL_M(\Omega)$ and $\Phi \in L_{\bar{M}}(\nu_n)$ such that
\begin{equation}\label{desi}
\iint_{\R^{2n}} (a(x,y,W)-\Phi)(W-D^s u)\,d\nu_n\geq 0,	
\end{equation}
for all $W\in L^{\infty}(\R^{2n},d\nu_n)$ with compact support then $a(x,y,D^su)=\Phi \in (W_0^sL_M(\Omega))^*$ that is
$$
\iint_{\R^{2n}}a(x,y,D^su)D^sv\,d\nu_n =\iint_{\R^{2n}}\Phi D^sv\,d\nu_n \qquad \forall v \in W_0^sL_M(\Omega).
$$
\end{lema}

\begin{teo}\label{pseudo}
Let $\Omega\subset \R^n$ be a bounded and open domain that satisfies the segment property. Then, the operator $\A$ is pseudomonotone. That is, If $\{u_i\}_{i\in \N}\subset \Dom(\A)$ is a sequence such that it fulfills
\begin{equation}\label{seq.pseudo}
\begin{cases}	
&u_i\to u\text{ for }\sigma(W_0^sL_M(\Omega),W^{-s}E_{\bar{M}}(\Omega))\\
&\A u_i\to f\in W^{-s}L_{\bar{M}}(\Omega)\text{ for }\sigma(W^{-s}L_{\bar{M}}(\Omega),W_0^sE_M(\Omega))\\
&\limsup_{i\to \infty}\langle \A u_i,u_i\rangle\leq \langle f,u\rangle
\end{cases}
\end{equation}
then 
\begin{equation}\label{tesis.pseudo}
\begin{cases}
&u\in \Dom(\A)
\\
&\A u=f
\\
&\langle \A u_i,u_i\rangle\to  \langle f,u\rangle\text{ if }i\to \infty.
\end{cases}
\end{equation}
\end{teo}

\begin{proof}[Proof of Theorem \ref{pseudo}]
Let $\{u_i\}_{i\in \N} \subset \Dom(\A)$ and $f\in  W^{-s}L_{\bar{M}}(\Omega)$ such that \eqref{seq.pseudo} is satisfied. 
We must prove that $u=\lim u_i$ and $f$ satisfy \eqref{tesis.pseudo}.


Let us first prove that the sequence $\{a(x,y,D^s u_i)\}_{i \in \N}$ remains bounded in $L_{\bar{M}}(\nu_n)$, indeed let $V \in L^{\infty}(\R^{2n},d\nu_n)$ with compact support, by the monotonicity condition \eqref{MC} we have 
\begin{equation}\label{mono}
\iint_{\R^{2n}} (a(x,y,D^s u_i)-a(x,y,V))(D^s u_i-V)\, d\nu_n\geq 0,
\end{equation}
which implies
\begin{align*}
\iint_{\R^{2n}} a(x,y,D^s u_i)V\,d\nu_n&\leq \iint_{\R^{2n}} a(x,y,D^s u_i)D^s u_i\,d\nu_n
\\
&-\iint_{\R^{2n}} a(x,y,V)D^s u_i\,d\nu_n
\\
&+\iint_{\R^{2n}} a(x,y,V)V\,d\nu_n
\\
&=I+II+III,
\end{align*}
the first integral remains bounded by \eqref{seq.pseudo}, the convergence $u_i\to u$ implies that the second integral is uniformly bounded in $i$ and the last integral is independent of $i$, therefore there exists a constant $C$, independent of $i$, such that
$$
\iint_{\R^{2n}} a(x,y,D^s u_i)V\,d\nu_n\leq C,
$$
for all $V\in L^{\infty}(\R^{2n},d\nu_n)$ with compact support and therefore for all $V\in E_M(\nu_n)$, then by the uniform boundedness principle we have $\|a(x,y,D^s u_i)\|_{L_{\bar{M}}(\nu_n)}\leq C$ uniformly in $i$ and so there exists a subsequence that we still denote by $\{a(x,y,D^s u_i)\}_{i\in \N}$ and   $\Phi\in L_{\bar{M}}(\nu_n)$ such that $a(x,y,D^su_i)\to \Phi$ in $\sigma(L_{\bar{M}}(\nu_n),E_M(\nu_n))$, using the above convergence and the fact $\A u_i \to f \in W^{-s}L_{\bar{M}}(\Omega)$ for $\sigma(W^{-s}L_{\bar{M}}(\Omega),W_0^sE_M(\Omega))$ we have
\begin{align*}
\langle f,v \rangle
&=\lim_{i\to\infty} \langle \A u_i,v\rangle
\\
&=\lim_{i\to\infty}\iint_{\R^{2n}}a(x,y,D^s u_i)D^sv\,d\nu_n
\\
&=\iint_{\R^{2n}}\Phi D^s v\,d\nu_n,
\end{align*}
for all $v\in W_0^sE_M(\Omega)$. This formula together with the density of $W_0^s
E_M(\Omega)$ in $W_0^s
L_M(\Omega)$ with respect to the $\sigma(L_M\times L_M(\nu_n),L_{\bar{M}}\times L_{\bar{M}}(\nu_n)$) topology allows us to extend $v$ to the space $W_0^sL_M(\Omega)$, that is
\begin{equation}\label{deff}
\langle f,v \rangle=\iint_{\R^{2n}}\Phi D^s v\,d\nu_n\qquad \forall v\in W_0^s L_M(\Omega).	
\end{equation}
Now taking limit in \eqref{mono} we obtain
$$
\iint_{\R^{2n}} (\Phi-a(x,y,V))(D^s u-V)\, d\nu_n\geq 0,
$$
for all $V\in L^{\infty}(\R^{2n},d\nu_n)$ with compact support and so by Lemma \ref{paraEulerLagrange} we have $a(x,y,D^s u)=\Phi$ in $(W_0^sL_M(\Omega))^*$, that is $\A u=f$.

On the other hand putting $u_i=u$ in \eqref{mono}, using $\A u=f$ and \eqref{deff} we have
\begin{align*}
\iint_{\R^{2n}} a(x,y,D^s u)V\,d\nu_n&\leq \iint_{\R^{2n}} \Phi D^s u\,d\nu_n
\\
&-\iint_{\R^{2n}} a(x,y,V)D^s u\,d\nu_n
\\
&+\iint_{\R^{2n}} a(x,y,V)V\,d\nu_n,
\\
&<\infty,
\end{align*}
for all $V\in E_M(\nu_n)$ that is $a(x,y,D^su)\in L_{\bar{M}}(\nu_n)$ and so $u\in \Dom(\A)$.

Finally, we prove that 
$$
\langle \A u_i, u_i \rangle \to \langle f, u \rangle \quad \text{for } i \to \infty.
$$
Let  
$$
L = \liminf_{i \to \infty} \iint_{\mathbb{R}^{2n}} a(x,y,D^s u_i) D^s u_i \, d\nu_n,
$$
we only need to prove that \( L \geq \langle f, u \rangle \). 

Using the monotonicity property again we have  
$$
\iint_{\mathbb{R}^{2n}} \left(a(x,y,D^s u_i) - a(x,y,D^{s,j}u) \right) (D^s u_i - D^{s,j}u) \, d\nu_n \geq 0,
$$
where $ D^{s,j}u = D^s u \chi_{R_j(u)}$. Taking the limit in \( i \) and rewriting, we obtain  
$$
L \geq \iint_{\mathbb{R}^{2n}} a(x,y,D^{s,j}u)(D^s u - D^{s,j}u) \, d\nu_n + \iint_{\mathbb{R}^{2n}} a(x,y,D^s u) D^{s,j}u \, d\nu_n.
$$
In $ R_j(u)$ the factor $D^s u - D^{s,j}u = 0$, and in $R_j(u)^c$ the factor $a(x,y,D^{s,j}u)=a(x,y,0)=0$, so the first integral in the above inequality is zero. Therefore  
$$
L \geq \iint_{R_j(u)}a(x,y,D^s u) D^s u \, d\nu_n,
$$
for arbitrary $j$, then  
$$
L \geq \iint_{\mathbb{R}^{2n}}a(x,y,D^s u) D^s u \, d\nu_n = \langle f, u \rangle.
$$
This concludes the proof of the theorem.
\end{proof}

\end{document}
